\subsection{The original key line and the shifted short family}
\label{sec:actual-key-short-span}

Put $g=f(p+q)=fr$.  On the original abelian family write
$\mathcal T=\pi_*\mathscr L$ for its rank-one theta-section line
and $\omega=\det\operatorname{Lie}(G)^*$.
The source's key-weight line is exactly
\[
                     \mathcal K=\mathcal T^8\otimes\omega^4.
\]
If the theta generator is multiplied by $a$ and the full tangent
frame is changed by $A$, the coefficient of a section of
$\mathcal K$ is multiplied by $a^{-8}(\det A)^4$.
This is the original homogeneous weight with parallel part zero
and theta tag $-4$.

\begin{proposition}\label{prop:actual-integral-key-shift}
On the actual good cover there is an algebraic generator $k$ of
the original generic key line, defined over the original coefficient
field, such that
\[
                  \operatorname{homN}_\sigma(k;z)=(2\pi)^{2g}
\]
at every original complex frame.  Its value at every original
integral eigenframe has nonnegative normalized valuation, including
at primes above $2$ and at ramified coefficient primes.  The section
and its generic inverse preserve the original cusp spaces under
their respective weight shifts.  In particular $k^m$ is an integral
twist of tag $-4m$, with constant norm $(2\pi)^{2gm}$, for every
integer $m\ge0$.
\end{proposition}

\begin{proof}
\emph{The all-prime key formula.}
The required classical theorem is the original integral key formula,
not just its formulation after inverting $2$.
Moret--Bailly works with the stack of principally polarized abelian
schemes with a symmetric theta divisor over $\mathbf Z$.
If $\mathrm{TH}=e^*\mathcal O(\Theta)$, his
$\Delta=\mathrm{TH}^{-2}\otimes\omega$ has a unit trivialization
$\varphi$ of $\Delta^4$ on this entire stack
\cite[\S0.1 and Theorem 0.2]{theta-key-moret-bailly}.
For our specified symmetric rigidified principal line,
$\mathscr L=\mathcal O(\Theta)\otimes\pi^*\mathrm{TH}^{-1}$
and $\pi_*\mathcal O(\Theta)=\mathcal O$ via its divisor section.
Thus $\mathcal T=\mathrm{TH}^{-1}$ and $\Delta^4=\mathcal K$.
Pull back $\varphi$ along this same original family.  This gives
$k$, and makes its coefficient a unit in any full integral Lie
basis and generating theta section whenever the family is abelian.
The assertion includes residue characteristic two.  A choice of the
permitted universal signs does not change any conclusion here;
the short-span consumer requires a generator of this original line,
not a further choice of theta-characteristic phase.

\emph{The exact metric.}
In Moret--Bailly's metric the divisor section has squared probability
$L^2$ norm $2^{-g/2}$, and the Hodge norm uses
$ (i/2)^g(-1)^{g(g-1)/2}\int\eta\wedge\bar\eta$.
His Theorem 3.3 gives $\|\varphi\|=(2\pi)^{4g}$, for both
theta parities.  Our theta line with its original rigidified metric
is $\mathrm{TH}^{-1}$ with the additional norm factor $2^{-g/4}$.
Our Hodge norm has the additional factor $\pi^{-g/2}$.
On the eighth and fourth powers, respectively, these give
\[
        2^{-2g}\pi^{-2g}(2\pi)^{4g}=(2\pi)^{2g}.
\]
This is exactly the source's normalized key norm; its Lie factor
uses the original $\xi$ and cancels in the displayed intrinsic
metrics\cite{theta-key-source}.

\emph{Extension through the original semiabelian charts.}
We justify the integral boundary assertion, which does not follow
from the key formula for abelian schemes alone.  Start with an
original integral test frame over a complete discrete valuation
ring $R$ whose generic family is abelian.  The actual DVR
uniformization gives its original Raynaud and period data.
After a finite extension, make the positive integral valuation
form of its periods divisible by $2M$ and write
\[
        \psi_y=\psi_y^0\pi^{Q(y)},\qquad
        \tau_{y,z}=\tau_{y,z}^0\pi^{b(y,z)},
        \quad Q(y)=b(y,y)/2,
\]
with all remaining factors unit generators.  The primitive identity
gives $b(y,w_0y)=Q(y)$, so it survives this division exactly.
These are the original arc data of
the established original arc theorem\cite{current-construction-record}.  Replace $\pi$ by $t$ and use
its constructed family over $A=R'[[t]]$.  Its abelian locus is
exactly $D(t)$, and its specialization at $t=\pi$ is the same
original family, line and level.

Use the generating theta section whose coefficient zero is the
chosen generator $\theta_B$ on the abelian quotient.  The formal
Raynaud isomorphism identifies the Lie algebra with that of the
constant Raynaud extension; choose a basis of this constant free
$R'$-module.  These give the actual generating theta section and
full Lie basis over $A$.  Replacing either by any other integral
generator or basis only multiplies the key coefficient by a unit.
In the chosen frames the pullback
of $k$ is a unit $F\in A[1/t]^\times$, by the all-prime abelian
key formula.  We show that its order at $t=0$ is zero.
Here is the precise coefficient-field reduction.  Put
$L=\operatorname{Frac}R'$.  The constant Raynaud abelian variety,
its rigidified line, a basis of its Lie algebra, $\theta_B$, the
point map on a basis of the finite-rank period lattice, and the
unit period factors on its generators descend to a finitely
generated subfield $L_0\subset L$ of characteristic zero.  Include
the original number fields and the finitely many level labels in
$L_0$.  The cube and bimultiplicative identities determine all
remaining period factors from these finite data.  After enlarging
$L_0$ by finitely many coefficients if necessary, all these
identities hold over $L_0$.  The same original arc construction over
$L_0[[t]]$ and its specified base change identify its theta section
and Lie frame with those just chosen over $L[[t]]$.  Thus the
Laurent coefficient is some $F_0\in L_0((t))$, whose image is $F$;
in particular $\operatorname{ord}_tF_0=\operatorname{ord}_tF$.

Choose an embedding $L_0\hookrightarrow\mathbf C$, extending a
chosen embedding of the original number field.  It is injective,
so the first nonzero Laurent coefficient remains nonzero and the
$t$-order is unchanged.  This does not embed the whole complete
valuation ring continuously in $\mathbf C$, nor specialize an
arbitrary formal series analytically.  The constant period data
and the monomials $t^{Q(y)},t^{b(y,z)}$ have an actual convergent
complex realization on a punctured disc.  The same toroidal chart
extends this realization over zero; completion recovers precisely
the arc by its proved uniqueness.  The algebraic key section is
a meromorphic section on that chart, so its pullback is a
meromorphic germ whose Laurent series is the image of $F_0$.
In this realization choose symplectic
coordinates adapted to the torus and the fixed abelian quotient.
The period matrix has the form
\[
        \Omega(t)=\Omega_0+
             \frac{\log t}{2\pi i}
                  \begin{pmatrix}B&0\\0&0\end{pmatrix},
                      \qquad B>0,
\]
up to fixed changes of these adapted coordinates.  The invariant
torus differentials are $dz_j/z_j$, and the remaining differentials
are lifts of a fixed basis on the abelian quotient.  They form a
frame extending across $t=0$.  Consequently the determinant of
the change from an actual generating Lie frame to this frame is
a holomorphic unit.

There is also no hidden power of $t$ in the theta scalar.  The
original expansion is indexed by the entire period lattice, has
factors $\psi_y^0t^{Q(y)}$, and has zero coefficient equal to
the nonzero generating theta section of the abelian quotient.
This is the exact zero-coefficient statement of
the established integral theta theorem\cite{current-construction-record}.  In adapted coordinates the standard
theta series has that same zero coefficient, up to a fixed nonzero
base-theta scalar.  Comparing the two coefficients shows that the
scalar $c(t)$ relating the actual generator to the standard theta
series is a holomorphic unit.  The positive definite $Q$ gives
normal convergence for sufficiently small $|t|$ on compact sets;
the nonzero coefficients are exponentially suppressed.  The
original integral recurrence, including $Q(0)=0$, is what rules
out a toric half-character shift or a fractional leading power.
Only the abelian quotient may retain a nonzero theta
characteristic.  This use of the actual primitive line is
essential: the key need not be a boundary unit for an arbitrary
symmetric principal line.

For the full tangent frame represented by
$(2\pi i)^{-1}M(t)$ in these theta coordinates, Gaussian
integration, with the same zero rigidification as the original
line, gives
\[
 \|\theta\|_{L^2}^2
       =|c(t)|^2\det(2\operatorname{Im}\Omega(t))^{-1/2},
 \qquad
 \|\eta\|^2
       =\frac{\det\operatorname{Im}\Omega(t)}
                    {\pi^g|(2\pi i)^{-g}\det M(t)|^2}.
\]
The second is exactly the source's dual Lie determinant norm.
There is no change of line metric or rigidifying scalar here.
Substitution into
$|F|\,\|\theta\|_{L^2}^{8}\|\eta\|^4=(2\pi)^{2g}$ gives
\[
                     |F(t)|=
                         \frac{|\det M(t)|^4}{|c(t)|^8}.
\]
The determinant and scalar are bounded above and away from zero
near $t=0$.  A nonzero meromorphic germ $t^e u(t)$, with $u(0)\ne0$,
has bounded modulus above and away from zero exactly when $e=0$.
Hence the image of $F_0$ has order zero, and so does the original
$F$.  This proves the no-pole assertion for both coefficient and
inverse, with no inference from a norm bound for just one of them.

The regular ring $R'[[t]]$ is normal.  A unit of $A[1/t]$ with
zero $t$-order is a unit of $A$, since all other height-one primes
already lie in $D(t)$.  In particular the prime $(\pi)$, including
the dyadic one, was checked by the all-prime abelian formula on
$D(t)$; the sole missing divisor was $(t)$, whose valuation is
exactly the generic $t$-order just computed.  Hence
$F\in A^\times$.  Specializing
$t=\pi$ proves that the key coefficient is an $R'$-unit on the
original degeneration.  Normalized valuations descend through
the finite extension, so it was an $R$-unit in a full integral
Lie basis.  This proves the boundary assertion without assuming
the source's key-unit-on-models declaration.

Finally an original integral eigenframe consists of actual vectors
of $\operatorname{Lie}(G/R)$.  Express their union in a full
integral basis, giving an integral matrix $U$, invertible over the
fraction field.  The key transformation law gives coefficient
$(\det U)^4u$, with $u\in R^\times$.  Its valuation is
nonnegative.  We have not replaced the ramified eigenspaces by a
free coefficient module or divided by their index.  The inverse
key section is a generic unit, with reciprocal constant complex
norm; the original holomorphy and cusp-vanishing conditions are
therefore preserved in both directions.  The stated power follows.
\end{proof}

\begin{theorem}\label{thm:actual-shifted-short-span}
Retain the hypotheses and notation of the actual short-theta
theorem, with $p,q\ge1$, $f>2$, $X=frL_b$, and
$\operatorname{fieldSz}(K)\le391X^2$.
Let the original weight $\lambda_{\tau,1}\ge\cdots\ge
\lambda_{\tau,r}$ satisfy
\[
 \lambda_{\tau,p}-\lambda_{\tau,p+1}\ge8fr+10,
 \qquad S=\sum_{\tau,j}|\lambda_{\tau,j}|,
 \qquad\log S\le10L_b.
\]
On every original component its cusp space of weight $\lambda$
and theta tag zero has a finite coefficient-field spanning family
in its original integral lattice, with every original Petersson
norm at most
\[
                  \exp\bigl(400000(X^3S+X^5)\bigr).
\]
This is field spanning by integral vectors; no integral module
basis or uniform rank bound is required.
\end{theorem}

\begin{proof}
Set $m=fr+1$, $n=8m$, and
$n_{0,\tau}=1+4m-\lambda_{\tau,p}$.
The shifted weight $\lambda'=\lambda+\operatorname{parWt}(n_0)$
has its first $p$ entries greater than $4m$ and its remaining
entries less than $-4m$.  Subtract $4m$ from the positive block
and subtract $4m$ from the negatives of the reversed negative
block.  This gives the original nonnegative antitone theta
exponents $a$ with $\operatorname{laOf}(a,n)=\lambda'$.
The associated size satisfies
\[
 \kappa_0\le(1+r)S+fr(4m+1)\le2XS+9X^2,
 \qquad
 \sum_\tau|n_{0,\tau}|\le S+f(4m+1)\le S+9X^2.
\]
These follow by the triangle inequality on the $r$ entries of
each half-embedding; the factor $1+r$ is retained.
Also $n\ge8fr+4$ and $\log\kappa_0\le20L_b$ under the displayed
assumptions: $X\le e^{2L_b}$ gives
$\kappa_0\le11e^{12L_b}\le e^{15L_b}$.
The actual short-theta theorem therefore supplies
an integral spanning family of weight $\lambda'$ and tag $4m$,
with log norm at most $30000X^2\kappa_0$.

Multiply this family by $k^m$ and the actual integral parallel
twist $t_{-n_0}$.  Their weights add to $\lambda$ and their tags
add to zero.  All three factors are integral at the same original
test frames.  The original common theta denominator remains the
same; no inverse key is used for the integral construction.
Multiplication by the product, whose generic inverse exists and
preserves cusp spaces, is a coefficient-field isomorphism of
the two cusp spaces.  The family still spans.

The norm bounds give
\[
 \begin{split}
 \log\|k^m t_{-n_0}\theta\|
 \le{}&30000X^2(2XS+9X^2)+2mfr\log(2\pi)\\
 &+\tfrac32rL_b(S+9X^2)
                 +\tfrac12\log\operatorname{rd}(K)+40fL_b+1.
 \end{split}
\]
Here $X\ge1$, $rL_b,fL_b\le X$,
$mfr\le2X^2$, $\log(2\pi)<2$, and
$\log\operatorname{rd}(K)\le391X^2$.
The right side is at most
$60001.5X^3S+270258X^5<400000(X^3S+X^5)$.
The same estimate holds for the original component Petersson
norm because each twisting factor has constant point norm.
The finite spanning family supplied above therefore has precisely
the asserted arithmetic and metric properties.
\end{proof}

\begin{corollary}\label{cor:actual-hs-short-families}
For the original asymptotic family, both the $h_0$ and
$h_{\rm sh}$ hypotheses of the Hilbert--Samuel degree estimate
hold on every original component, with its stated integral lattice
and Petersson norms.  One may take the same parameter threshold
used by the original construction, increased so that
$L_b\ge2400000$; no auxiliary field or new normalization is needed.
\end{corollary}

\begin{proof}
We spell out the two comparisons, retaining the exact original
weights.  Write $D=f\binom r2$ for the flag dimension,
$\lambda_i=P.\operatorname{wt}(i)$, $0\le i\le T$, and
$S_i=\sum_{\tau,j}|\lambda_{i,\tau,j}|$.
The original numerical parameter bounds are
\[
 \begin{gathered}
 f\ge8,\quad r\ge100,\quad \kappa\ge1,\quad
 \log\kappa\le2L_b,\quad PP\ge X^{12},\quad X=frL_b,\\
 S_i\le X^2\kappa\quad(0\le i\le T),\qquad
 \sum_{\tau,j}|u_{a,\tau,j}|\le
         X^2(\kappa/T+1)\quad(a=1,2).
 \end{gathered}
\]
These are the original size inequalities, not an extra estimate
on the arithmetic fields or cusp rank.  Their exact supplier and
consumer ranges are recorded in the source certificate
\cite{theta-key-source}.

The original weight formula is
$\lambda_i=\kappa_{\rm flag}+u_2+(D+i)u_1$.
The canonical flag weight has consecutive gaps at least $-2$;
the defining remainder inequalities and the original large-gap
bound on $u_1$ give consecutive gaps at least $D+1$ for every
$\lambda_i$.  Since
$D+1\ge8fr+10$ for the displayed $f,r$, each $\lambda_i$ is
antitone and has the exact central gap required above.  Moreover
\[
 \log S_i\le2\log(fr)+2\log L_b+\log\kappa\le6L_b<10L_b.
\]
The theta-cover hypotheses and the coefficient-field size bound
are the same as in the preceding theorem, by the original cover
laws.  Thus it applies simultaneously, with $C=400000$, to all
these weights and every component.

Put $Y_0=\kappa/T+1\ge1$.  The flag weight has total absolute
size at most $fr^2\le X^2$.  Hence the exact formula at zero and
the two bounds on $u_a$ imply
\[
             S_0\le(2+D)X^2Y_0\le3DX^2Y_0.
\]
With $X\ge L_b\ge6C$, the bound from the theorem is at most
\[
 C(X^3S_0+X^5)
       \le4CDY_0X^5\le D\,PP\,Y_0.
\]
The final inequality follows from $4C\le X^7$ and $PP\ge X^{12}$.
This gives precisely $h_0$.

For every $0\le i\le T$, the same theorem gives
\[
 C(X^3S_i+X^5)
       \le CX^5(\kappa+1)
       \le2CX^5\kappa\le PP\,\kappa,
\]
using $\kappa\ge1$, $2C\le X$, and $X^6\le X^{12}\le PP$.
This gives precisely $h_{\rm sh}$, including all its indices and
all embeddings.  The integrality and coefficient-field spanning
were already supplied by the constructed families, so no
integral module-basis conversion, rank estimate, or covolume
argument intervenes.  Finally the original host parameter is
$L_b=x^\varepsilon$ with $\varepsilon>0$; the added constant
threshold is obtained by increasing the original starting $x$,
as its quantified asymptotic statement permits.
\end{proof}
